\section{Desarrollo}

\subsection{Primera optimización: Evitar recálculos}

Se agrega una lista con la suma de los divisores de los números que se van iterando y luego se utiliza para hacer comparaciones y evitar 
recálculos de sumas de los mismos valores. Esta optimización logró reducir el tiempo de ejecución de 330 a 205 segundos para 100\,000 valores.

\textbf{Implementacion:} 

\begin{lstlisting}[language=Python]

    import time
    def amigos(MAX):
        sumas = []
        t1=time.time()
        for i in range(MAX):
            s=0
            for j in range (1,i-1+1):
                if i%j==0:
                    s+=j
            sumas.append(s)
            if s > i:
                continue
            else:
                s2 = sumas[s]
                if i == s2:
                    print(i,s)       
        t2=time.time()    
        print(t2-t1)

    amigos(100000)
\end{lstlisting}

\subsection{Segunda optimización: búsqueda de divisores}

Dada que la actual búsqueda de divisores para cada i es de O(i) y hay n valores, esto concluye en que para este paso
la complejidad temporal es O(n * n).

Esta optimización busca hacerse acotando el rango de búsqueda para cada i, reemplazando las n probabilidades por $\sqrt{n} + 1$
ya que si un número tiene un divisor mayor que su raíz cuadrada, necesariamente tiene un divisor menor que su raíz cuadrada.

\textbf{Justificación matemática:} 

Para entender por qué podemos acotar la búsqueda de divisores hasta $\sqrt{n}$, consideremos la siguiente demostración:

\begin{itemize}
    \item Sea $n \in \mathbb{N}$ un número entero positivo
    \item Sea $d$ un divisor de $n$ tal que $d > \sqrt{n}$
    \item Por definición de divisor: $\exists \, d' \in \mathbb{N}$ tal que $n = d \cdot d'$
\end{itemize}

\textbf{Demostración:}
\begin{align}
    d > \sqrt{n} &\implies \frac{n}{d} < \frac{n}{\sqrt{n}} \\
    &\implies d' = \frac{n}{d} < \sqrt{n}
\end{align}

Por lo tanto, por cada divisor $d > \sqrt{n}$ existe necesariamente un divisor complementario $d' < \sqrt{n}$. 

\textbf{Conclusión:} Solo necesitamos iterar hasta $\lfloor\sqrt{n}\rfloor$ para encontrar todos los pares de divisores. Cuando encontramos un divisor $i$, añadimos tanto $i$ como su complementario $\frac{n}{i}$ a la suma total.

Se implementa una función \texttt{suma\_de\_divisores} que utiliza esta propiedad, iterando solo hasta $\sqrt{n}$ y añadiendo ambos divisores cuando encuentra un par. Esto reduce la complejidad temporal de O(n) a O($\sqrt{n}$) por cada número.

\textbf{Implementacion:} 
\begin{lstlisting}[language=Python]
import math

def suma_de_divisores(numero):
    if numero <= 1:
        return 0
    suma = 1  # 1 siempre es divisor
    sqrt_n = int(math.sqrt(numero))
    
    for i in range(2, sqrt_n + 1):
        if numero % i == 0:
            suma += i
            if i != numero // i:  # Evitar contar dos veces el mismo divisor
                suma += numero // i
    
    return suma

def amigos(MAX):
    sumas = [0] * MAX
    t1 = time.time()
    for i in range(MAX):
        sumas[i] = suma_de_divisores(i)
        s = sumas[i]
        if s > i:
            continue
        else:
            s2 = sumas[s]
            if i == s2:
                print(i, s)       
    t2 = time.time()    
    print(t2 - t1)
\end{lstlisting}